\chapter{Governance, Purview, and GDPR Erasure}
\index{Purview}\index{governance}\index{sensitivity labels}\index{lineage}\index{endorsement}\index{GDPR}\index{right to erasure}\index{crypto-shredding}\index{data mesh}\index{legal hold}%

\begin{objectives}
\begin{itemize}
    \item Use the Purview hub and tenant administration surfaces for Fabric governance.
    \item Apply sensitivity labels to Fabric items and reason about downstream inheritance and legal-hold interaction.
    \item Trace end-to-end lineage and use impact analysis before making changes.
    \item Execute a GDPR right-to-erasure request across Delta tables, history files, KQL retention, caches, and backups.
    \item Implement crypto-shredding with per-subject keys in Azure Key Vault, and contrast it with row-level \texttt{DELETE}.
    \item Architect a five-domain data mesh with centralized Purview governance and per-domain erasure runbooks.
\end{itemize}
\end{objectives}

\begin{crosscloud}
Metadata governance---catalog, classification, lineage, and compliance workflows---is a distinct platform layer in every cloud, increasingly unified across the estate rather than per-service.

\vspace{2mm}
{\footnotesize
\begin{tabular}{@{}p{2.8cm}p{3.2cm}p{3.2cm}p{3.2cm}p{3.2cm}@{}}
\toprule
\textbf{Capability} & \textbf{Fabric} & \textbf{AWS} & \textbf{GCP} & \textbf{Snowflake} \\
\midrule
Catalog/governance & Microsoft Purview & DataZone / Glue Catalog + Macie & Dataplex & Horizon Catalog \\
Sensitivity labels & Microsoft Information Protection labels on items & Macie classifications & Sensitive Data Protection / policy tags & Object tagging, masking policies \\
Lineage & Fabric lineage view + Purview & Glue/DataZone lineage & Dataplex lineage & Access history / lineage \\
Endorsement & Promoted / Certified & Business glossary approvals & Data product annotations & Listing certification \\
Erasure mechanics & Crypto-shredding, Delta VACUUM coordination & S3 object deletion + Glacier waits & TTLs, crypto-shredding (CMEK) & Time Travel/Fail-safe windows \\
Legal hold & Purview eDiscovery/holds & Legal holds (S3 Object Lock) & Vault holds & Retention + legal processes \\
\bottomrule
\end{tabular}}

\vspace{1mm}
Databricks governs through Unity Catalog with Purview or Unity lineage; MongoDB through Atlas auditing and field-level encryption. The hard, portable lesson: \emph{erasure is a propagation problem across caches, history, and backups---row deletes alone never satisfy it}.
\end{crosscloud}

\begin{keyterms}
\textbf{Purview hub}: Fabric's integrated governance surface for insights, sensitivity, and lineage. \quad
\textbf{Sensitivity label}: a classification (for example \emph{Highly Confidential}) applied to an item, traveling downstream. \quad
\textbf{Endorsement}: Promoted (vouched) or Certified (officially approved) status on a data product. \quad
\textbf{Lineage}: the item-to-item graph of data flow used for impact analysis. \quad
\textbf{Crypto-shredding}: erasing a data subject by destroying their encryption key, rendering every encrypted copy unreadable. \quad
\textbf{Legal hold}: a preservation obligation that overrides erasure until released, recorded as a documented exception.
\end{keyterms}

\section{Week 22 Roadmap}
Week~22 governs the estate. Monday covers the Purview hub and tenant administration. Tuesday covers sensitivity labels and their downstream inheritance---including the legal-hold interaction. Wednesday is the heavyweight: lineage, endorsement, impact analysis, and the GDPR erasure propagation path, comparing crypto-shredding against row-level \texttt{DELETE}. Thursday is the hands-on project: label a lakehouse, trace lineage to a Power BI report, and execute a crypto-shredding erasure. Friday architects the five-domain data mesh with centralized governance.

Everything before this week built capability; this week builds accountability. Regulators, auditors, and---most importantly---data subjects do not care how elegant the pipelines are; they care whether the organization can prove what it holds and delete what it must \cite{microsoftPurviewFabric,microsoftFabricLineage}.

\begin{notebox}
The week's hard truth, stated early: naive \texttt{DELETE} is insufficient for GDPR erasure in Fabric. Delta time travel keeps deleted rows readable in history until \texttt{VACUUM} expires them; KQL data persists per retention; semantic-model caches hold copies; clones and backups re-create erased rows on restore. Erasure is a cross-layer propagation design, not a SQL statement.
\end{notebox}

\section{Monday: Purview Hub and Tenant Administration}
\index{Purview hub}\index{tenant administration}
The Purview hub inside Fabric surfaces governance posture to those who own it: sensitivity-label coverage, endorsement distribution, item inventory, and usage insights across the tenant \cite{microsoftPurviewFabric}. Tenant administration (the admin portal) sets the guardrails within which all workspaces operate: which features are enabled, which capacities exist, which export and sharing behaviors are permitted \cite{microsoftFabricTenantSettings}.

\begin{definitionbox}
The \textbf{Purview hub} is Fabric's integrated governance surface: it aggregates label coverage, endorsement, inventory, and usage signals so governance is observed continuously rather than audited annually.
\end{definitionbox}

The data engineer's Monday skill is reading the hub as a work queue: unlabeled items containing personal data, unendorsed tables feeding executive reports, domains without owners. Governance at scale is not a policy PDF; it is these queues trending toward zero.

\section{Tuesday: Sensitivity Labels and Inheritance}
\index{sensitivity labels}\index{label inheritance}\index{legal hold}
Sensitivity labels from Microsoft Information Protection apply to Fabric items---lakehouses, warehouses, semantic models, reports---and classify the data they contain \cite{microsoftPurviewFabric}. Two behaviors matter most:

\textbf{Downstream inheritance.} When a labeled item feeds another---a labeled lakehouse feeding a semantic model feeding a report---the label flows downstream, so classification follows the data through the serving layers. Inheritance is a default to be audited, not a guarantee to be assumed: verify in the hub that every downstream item of a \emph{Highly Confidential} source carries the label.

\textbf{Legal-hold interaction.} A retention or legal-hold policy overrides an erasure request: data under hold must not be deleted, even for a valid GDPR request. The resolution is procedural, not technical---the hold is recorded as a documented exception in the audit register, the erasure executes when the hold releases, and the register proves both the obligation and its handling \cite{microsoftPurviewFabric}.

\begin{pitfall}
Labels classify; they do not by themselves enforce. A \emph{Highly Confidential} label without matching access control (Chapter~21) and without export restrictions is a red sticker on an open door. Label, then verify the enforcement chain behind the label.
\end{pitfall}

\section{Wednesday: Lineage, Endorsement, and the Erasure Path}
\index{lineage}\index{endorsement}\index{GDPR!erasure path}\index{crypto-shredding}
\subsection{Lineage and Impact Analysis}
Fabric records item-to-item lineage: this pipeline feeds that lakehouse, which feeds that semantic model, which feeds that report \cite{microsoftFabricLineage}. Lineage's daily value is \emph{impact analysis}: before changing a table's schema, renaming a column (Chapter~9's contract break), or retiring an item, the lineage view lists every downstream consumer who will feel it. Endorsement completes the signal: \emph{Promoted} means the owner vouches; \emph{Certified} means the organization certifies---and certification is what tells an analyst which of the fourteen \texttt{sales} tables is the real one.

\subsection{The GDPR Erasure Propagation Path}
A right-to-erasure request must reach every copy \cite{microsoftPurviewFabric}:

\begin{enumerate}
    \item OneLake raw zones (\texttt{Files} folders, landed extracts).
    \item Bronze/Silver/Gold Delta tables---\emph{including time-travel history files}, which keep deleted rows readable until \texttt{VACUUM} expires them.
    \item Eventstream/KQL retention windows---telemetry ages out per policy, not per request.
    \item Downstream derived tables and aggregates containing the subject's data.
    \item Semantic-model caches holding rows in memory.
    \item Clones and backups, which re-create erased rows on restore.
\end{enumerate}

\subsection{Crypto-Shredding versus Row-Level DELETE}
\textbf{Row-level DELETE} removes current rows but leaves history, caches, and backups readable---erasure becomes a multi-week coordination of \texttt{VACUUM} schedules, retention windows, and cache clears, each with its own failure mode. \textbf{Crypto-shredding} encrypts each subject's PII columns with a per-subject key in Azure Key Vault; deleting the key renders every copy---current, historical, cached, backed-up---permanently unreadable in one auditable act \cite{microsoftFabricSecurity}. The trade-off: crypto-shredding must be designed in from the start (you cannot shred what you never encrypted per subject), and it makes the key vault a compliance-critical system.

\begin{table}[H]
\centering
\small
\begin{tabular}{@{}p{3.6cm}p{5.0cm}p{5.0cm}@{}}
\toprule
\textbf{Property} & \textbf{Row-level DELETE} & \textbf{Crypto-shredding} \\
\midrule
History files & Readable until VACUUM & Unreadable instantly \\
Caches/backups & Persist until cleared/rotated & Unreadable instantly \\
Design cost & Low, retrofit-able & High, must be designed in \\
Auditability & Multi-step evidence chain & One key-deletion event \\
Key vault criticality & None & Compliance-critical system \\
\bottomrule
\end{tabular}
\caption{Two erasure strategies, honestly compared.}
\label{tab:chapter22-erasure}
\end{table}

\begin{pitfall}
Crypto-shredding without key-vault hygiene is a single point of catastrophic data loss: accidental or malicious bulk key deletion is unrecoverable by design. Key lifecycle---creation, rotation policy, deletion approval workflow, soft-delete windows---is the governance object that makes shredding safe.
\end{pitfall}

\section{Thursday: Hands-on Project}
\index{hands-on lab}\index{crypto-shredding!hands-on}
The Week~22 lab applies a \emph{Highly Confidential} label to a lakehouse, traces lineage from lakehouse through pipeline to Power BI report, and then implements crypto-shredding: per-subject keys in Key Vault, PII encrypted in a notebook, and an erasure request executed by deleting one subject's key---verifying the data is unrecoverable in current tables and Delta history alike, while a row-level \texttt{DELETE} leaves rows readable via time travel until \texttt{VACUUM}.

\subsection{Goal and Architecture}
Workspace \texttt{fabric-de-week22} with lakehouse \texttt{lh\_pii} containing \texttt{customers\_pii}, a pipeline and report consuming it, and an Azure Key Vault. The lab proves both erasure strategies and their asymmetry.

\subsection{Step-by-Step Actions}
\begin{enumerate}
    \item Apply the \emph{Highly Confidential} label to \texttt{lh\_pii}; verify inheritance to the downstream model and report in the Purview hub.
    \item Open the lineage view; capture the path lakehouse $\rightarrow$ pipeline $\rightarrow$ model $\rightarrow$ report.
    \item Generate per-subject AES keys in Key Vault for a small synthetic subject set.
    \item Encrypt the PII columns per subject in a notebook; store ciphertext in the table, keys only in the vault.
    \item Erasure by \texttt{DELETE}: remove subject A's rows; prove they remain readable via \texttt{VERSION AS OF}.
    \item Erasure by shredding: delete subject B's key; prove their ciphertext is undecryptable in current tables \emph{and} history.
    \item Record both operations in an audit register, including any legal-hold check performed first.
\end{enumerate}

\begin{lstlisting}[language=Python,caption={Per-subject encryption pattern (key resolution from Key Vault only).},label={lst:chapter22-shred}]
from pyspark.sql import functions as F
from cryptography.fernet import Fernet

# Keys live ONLY in Azure Key Vault; resolved per subject at run time.
def encrypt_with_subject_key(subject_id: str, plaintext: str) -> str:
    key = keyvault.get_secret(f"pii-key-{subject_id}")   # vault reference, never literal
    return Fernet(key.encode()).encrypt(plaintext.encode()).decode()

encrypt_udf = F.udf(encrypt_with_subject_key, "string")

df = spark.read.table("customers_pii_raw")
(df.withColumn("email_enc", encrypt_udf("customer_id", "email"))
   .drop("email")
   .write.format("delta").mode("overwrite")
   .saveAsTable("customers_pii"))

# Erasure = delete the key in Key Vault: every copy, including
# Delta history and backups, becomes permanently undecryptable.
\end{lstlisting}

\subsection{Validation Checklist}
\begin{itemize}
    \item The label appears on the lakehouse and inherited downstream items.
    \item The lineage graph is captured and matches the intended flow.
    \item \texttt{DELETE}+time-travel demonstrates the residual-history problem.
    \item Key deletion demonstrates complete, instant erasure including history.
    \item Keys never appear in code, logs, or tables---only vault references.
    \item The audit register records both operations and the legal-hold check.
\end{itemize}

\section{Friday: Use Case and Architecture}
\index{data mesh}\index{governance!federated}
A global conglomerate operates five independent business domains---Retail, Logistics, Finance, Health, Media---each owning its data products in separate Fabric workspaces. The mandate: data-mesh autonomy with centralized governance through Purview.

\subsection{The Mesh Design}
\begin{itemize}
    \item \textbf{Domains as Fabric domains}: each business unit's workspaces grouped under a domain with a named owner---autonomy with accountability \cite{microsoftDomains}.
    \item \textbf{Federated computation, centralized policy}: domains build their own pipelines and products; Purview supplies the shared catalog, label taxonomy, lineage, and endorsement standards.
    \item \textbf{Contracts between domains}: cross-domain data sharing happens through endorsed data products with Week~9-style contracts, never through ad hoc lakehouse grants.
    \item \textbf{The per-domain erasure runbook}: each domain implements the Wednesday propagation path for its own Bronze/Silver/Gold lakehouses and KQL databases---a subject request fans out to five domains, and each domain's runbook executes the same steps against its own estate, with the central register collecting completion evidence.
\end{itemize}

\begin{figure}[H]
    \centering
    \begin{tikzpicture}[node distance=7mm and 8mm]
        \node[stage, minimum width=2.2cm] (d1) at (0,2.2) {Retail\\domain};
        \node[stage, minimum width=2.2cm] (d2) at (0,0) {Logistics\\domain};
        \node[stage, minimum width=2.2cm] (d3) at (0,-2.2) {Finance\\+ 2 more};
        \node[stage, fill=orange!10, minimum width=3.0cm] (pur) at (5.4,0) {Microsoft Purview\\catalog, labels,\\lineage, register};
        \node[doc, minimum width=2.8cm, minimum height=1.2cm] (dpo) at (10.4,0) {Erasure\\request\\(DPO)};

        \draw[biarrow] (d1) -- (pur);
        \draw[biarrow] (d2) -- (pur);
        \draw[biarrow] (d3) -- (pur);
        \draw[arrow] (dpo) -- (pur);

        \node[note, text width=11.5cm, align=center] at (5.2,-3.4) {Central catalog and register; per-domain runbooks execute the propagation path locally.};
    \end{tikzpicture}
    \caption{Five-domain data mesh with centralized Purview governance.}
    \label{fig:chapter22-mesh}
\end{figure}

\begin{pitfall}
Centralized governance with centralized execution recreates the bottleneck the mesh was designed to break. Purview centralizes \emph{policy and visibility}; domains keep \emph{execution}. The erasure runbook is the test case: centrally defined, domain-executed, centrally evidenced.
\end{pitfall}

\section{Operational Guidance for Week 22}
Governance fails at the edges: unlabeled new items, unowned domains, untested erasure. Schedule the hub review, the inheritance audit, and---annually at minimum---an erasure fire drill on a synthetic subject.

\subsection{Performance and Reliability}
Lineage and labels add negligible runtime cost; their cost is discipline. Erasure coordination with \texttt{VACUUM} schedules must respect Chapter~4's retention rules---shorter retention eases erasure and complicates recovery; the balance is a documented decision.

\subsection{Security and Governance Checklist}
\begin{itemize}
    \item Label coverage monitored in the hub; inheritance audited, not assumed.
    \item Legal holds recorded in the register and checked before any erasure.
    \item Lineage reviewed before any schema-breaking change.
    \item Certification reserved for products passing the published criteria.
    \item Erasure drill executed annually with per-domain evidence.
    \item Key-vault lifecycle (approvals, soft-delete, rotation) documented and tested.
\end{itemize}

\section{Interview Q\&A}
\begin{enumerate}
    \item \textbf{Q: What does a sensitivity label pass on to downstream items?}\\
    A: The classification itself---labels inherit downstream through lineage, so a \emph{Highly Confidential} lakehouse's models and reports carry the label; inheritance is audited, not assumed.
    \item \textbf{Q: What is the difference between Promoted and Certified endorsement?}\\
    A: Promoted is the owner's vouch; Certified is the organization's official approval against published criteria---the signal that tells consumers which table is the real one.
    \item \textbf{Q: How does lineage help with impact analysis before a change?}\\
    A: It enumerates every downstream consumer of an item, so a schema change or retirement is reviewed against the actual dependency list rather than discovered by broken reports.
    \item \textbf{Q: Why is row-level DELETE insufficient for GDPR erasure in Fabric?}\\
    A: Delta history keeps rows readable until \texttt{VACUUM}, KQL retains per policy, caches hold copies, and backups restore them; DELETE addresses only the current table state of one layer.
    \item \textbf{Q: How does crypto-shredding sidestep all of those?}\\
    A: PII is encrypted per subject with keys in Key Vault; deleting the key renders every copy---current, history, cache, backup---unreadable in one auditable act, with no dependence on per-layer cleanup timing.
    \item \textbf{Q: What is a legal hold's interaction with erasure?}\\
    A: A hold overrides the erasure obligation until release; it is recorded as a documented exception in the audit register, and erasure executes on release---procedural resolution, not technical.
\end{enumerate}

\section{Further Reading}
Review Microsoft Purview and Fabric governance \cite{microsoftPurviewFabric}, lineage \cite{microsoftFabricLineage}, sensitivity labels \cite{microsoftFabricLabels}, and domains \cite{microsoftDomains}. The Delta retention interactions come from Chapter~4 \cite{microsoftFabricVacuum,deltaLakeDocs}; access layers from Chapter~21 \cite{microsoftFabricSecurity}.

\section*{Key Takeaways}
\begin{itemize}
    \item The Purview hub turns governance into observable queues: label coverage, endorsement, ownership.
    \item Labels classify and inherit; enforcement is a separate, verified chain.
    \item Lineage is the impact-analysis tool: read it before breaking anything.
    \item Erasure is propagation across raw zones, Delta history, KQL retention, caches, and backups---DELETE alone never suffices.
    \item Crypto-shredding makes erasure one auditable key-deletion, at the price of designed-in encryption and key-vault criticality.
    \item Legal holds override erasure and live in the audit register as documented exceptions.
    \item Data mesh centralizes policy and visibility through Purview while domains execute---including their own erasure runbooks.
\end{itemize}

\section{Exercises}
\begin{enumerate}
    \item \textbf{(Conceptual)} Explain why a label without enforcement is worse than no label for audit confidence.
    \item \textbf{(Conceptual)} Argue both sides: when is row-level DELETE an acceptable erasure strategy?
    \item \textbf{(Hands-on)} Extend the lab: execute \texttt{VACUUM} after the DELETE and show time-travel readability ends.
    \item \textbf{(Hands-on)} Write the per-domain erasure checklist for one domain, timed step by step.
    \item \textbf{(Design)} Design the central erasure register: schema, evidence fields, and the query proving a request is fully satisfied.
    \item \textbf{(Design)} Draft the key-deletion approval workflow that makes crypto-shredding safe against accidents.
    \item \textbf{(Interview)} In five sentences, walk a DPO through what happens the day an erasure request arrives.
    \item \textbf{(Operational)} Write the annual erasure fire-drill procedure with pass/fail criteria.
\end{enumerate}

\section{Capstone Project: Auditable Erasure Pipeline}\label{sec:ch22-capstone}
\index{capstone project}\index{GDPR!capstone}\index{crypto-shredding!capstone}

\begin{capstonebox}
\textbf{Mission.} Build the complete erasure capability: labeled and lineage-traced PII estate, per-subject encryption, both erasure strategies demonstrated with their asymmetry proven, a legal-hold check, and a central audit register producing completion evidence for a sample request across two ``domains.''

\textbf{Estimated time.} 6--8 hours.

\textbf{What you'll practice.}
\begin{itemize}
    \item Label application, inheritance verification, and lineage capture.
    \item Per-subject encryption with vault-only key resolution.
    \item Executing and evidencing both erasure strategies.
    \item The register schema and completion query that satisfy an auditor.
\end{itemize}
\end{capstonebox}

\subsection*{Scenario}
Contoso Group's Data Protection Officer requires a working demonstration---not a policy deck---that a subject erasure request propagates correctly through the Fabric estate, honors legal holds, and produces auditable evidence.

\subsection*{Requirements}
\begin{itemize}
    \item Two workspace ``domains,'' each with a PII lakehouse, labels applied and inheritance verified.
    \item Per-subject encryption with keys in Key Vault; no key material in code or logs.
    \item The propagation checklist executed per domain: raw, tables, history, derived, caches.
    \item The DELETE-versus-shredding asymmetry demonstrated with time-travel evidence.
    \item The legal-hold check integrated into the workflow.
    \item The audit register with a completion query across both domains.
\end{itemize}

\subsection*{Implementation Steps}
\begin{enumerate}
    \item Stand up both domains' estates with labels and lineage captured.
    \item Implement per-subject encryption.
    \item Execute the DELETE path with time-travel and \texttt{VACUUM} evidence.
    \item Execute the shredding path with history-included unreadability evidence.
    \item Run one full request through the register, including the hold check.
\end{enumerate}

\subsection*{Deliverables}
\begin{itemize}
    \item Label and lineage evidence per domain.
    \item The encryption notebook and vault configuration (references only).
    \item Both strategies' evidence captures.
    \item The register schema, rows, and completion query output.
    \item The per-domain runbook documents.
\end{itemize}

\subsection*{Acceptance Criteria}
\begin{itemize}
    \item[\(\square\)] Labels inherit downstream in both domains, verified in the hub.
    \item[\(\square\)] DELETE leaves history readable until \texttt{VACUUM}, demonstrated.
    \item[\(\square\)] Key deletion renders the subject unreadable everywhere including history, demonstrated.
    \item[\(\square\)] A held subject is skipped and the exception recorded.
    \item[\(\square\)] The register's completion query returns full satisfaction for the executed request.
    \item[\(\square\)] No keys, secrets, or real personal data appear in any artifact.
\end{itemize}

\subsection*{Stretch Goals}
\begin{itemize}
    \item Add the third domain as a KQL/telemetry estate and extend the checklist to retention windows.
    \item Automate register fan-out: request in, per-domain tasks created, evidence collected.
    \item Add the key-deletion approval workflow with dual sign-off.
    \item Produce the annual fire-drill report template from the register.
\end{itemize}
